% =============================================================================
\section{Comentário aos resultados}
\label{sec:comentario}
% =============================================================================

\subsection*{\textcolor{bom}{Pontos bons}}

\begin{itemize}[itemsep=3pt]
  \item \textbf{Cinemática direta validada de quatro formas}: pose de
        referência, segunda tabela do fabricante, modelo 3D do robô e
        controlador.
  \item \textbf{Trajetória com folga em todos os limites} (no máximo
        \SI{44}{\percent} de um limite) e \textbf{sem contacto}, verificado
        em simulação com o volume real do robô antes de o mover.
  \item \textbf{Seguimento previsto e medido coincidem}: o erro é um ciclo
        de atraso (\SI{1.1}{\milli\meter} em média à velocidade planeada),
        reprodutível, até \SI{43}{\percent} acima da velocidade planeada.
  \item \textbf{As dúvidas da primeira versão foram resolvidas com dados}: a
        configuração lida do controlador refutou as hipóteses de um limite de
        velocidade e de carga mal configurada, e identificou a calibração como
        origem da diferença para a nossa cinemática.
\end{itemize}

\subsection*{\textcolor{mau}{O que não correu tão bem}}

\begin{itemize}[itemsep=3pt]
  \item \textbf{A primeira versão, no espaço cartesiano, não era seguida à
        velocidade planeada} (erro médio de \SI{8}{\milli\meter}), apesar de
        respeitar os limites publicados; só o redesenho nas juntas o resolveu,
        e a causa exata não foi isolada.
  \item \textbf{Sem medição externa.} O erro de seguimento é relativo à pose
        que o controlador reporta; a posição real da ponta só é conhecida a
        menos da calibração e da repetibilidade do robô.
  \item \textbf{O swing é lento} face ao do Lab~1
        (\SI{0.22}{\meter\per\second} no impacto), consequência das dimensões
        do robô e do limite de \SI{500}{\milli\meter\per\second}.
\end{itemize}

\subsection*{\textcolor{azul}{O que podíamos melhorar}}

\begin{itemize}[itemsep=3pt]
  \item \textbf{Medir a ponta com um instrumento externo}, para separar o
        erro de seguimento do erro do modelo.
  \item \textbf{Calibrar a cinemática} com poses medidas em todo o espaço de
        trabalho.
  \item \textbf{Aproximar a velocidade do limite do robô}: a trajetória foi
        seguida sem degradação a \SI{306}{\milli\meter\per\second}, pelo que
        há margem para um downswing mais curto.
\end{itemize}

\newpage
% =============================================================================
\section*{Contribuição dos membros}
Todos os membros contribuíram de forma igual para a realização do projeto.
Cada um ficou responsável por uma área, sem deixar de participar nas restantes:
\begin{itemize}
  \item \textbf{Afonso Cruz:} modelação: cinemática direta e modelo 3D com
        verificação de contacto.
  \item \textbf{Eduardo Batista:} planeamento: trajetória de swing e limites.
  \item \textbf{Henrique Póvoa:} execução: scripts do robô e comparação com a
        inversa do controlador.
\end{itemize}
As sessões de laboratório e a redação do relatório foram feitas por todos.

\section*{Declaração de uso de IA}
Foi usada uma LLM como apoio na geração de código Python e
na revisão do texto. Todos os membros compreendem integralmente o trabalho
entregue.
